% Ngân hàng bài dummy (tự sinh để test): Phương trình bậc hai một ẩn và định lý Viète (lớp 9).
% Quy ước:
%   - Mỗi bài: \begin{problem}[id] ... \end{problem}; lời giải lồng bên trong.
%   - id: duy nhất, không dấu, không đổi, không dùng lại. Dạng <lớp><phân môn>-<chủ đề>-<số>.
%   - Mỗi bước lời giải là 1 đoạn, các bước cách nhau 1 dòng trống; mỗi bước có lý do.
%   - Dòng cuối lời giải: Đáp số: $\boxed{...}$
%   - Cách giải khác: thêm \begin{solution}[cách 2].
%   - Metadata bằng comment ngay sau \begin{problem}: % nguon:, % lop:, % kc_du_kien:, % do_kho_du_kien:
\documentclass[12pt,a4paper]{article}
\usepackage[utf8]{vietnam}
\usepackage{amsmath,amssymb}
\newenvironment{problem}[1][]{\par\medskip\noindent\textbf{Bài #1.}\ }{\par}
\newenvironment{solution}[1][]{\par\noindent\textit{Lời giải #1.}\par}{\par}

\begin{document}

% ---------------------------------------------------------------- phương trình khuyết
\begin{problem}[9DS-PT2-001]
% nguon: dummy (tự sinh để test)
% lop: 9
% kc_du_kien: 9DS-PT2-01
% do_kho_du_kien: 1
Giải các phương trình sau:

a) $3x^2 - 12x = 0$;

b) $2x^2 - 18 = 0$;

c) $5x^2 + 20 = 0$.

\begin{solution}
a) Phương trình khuyết hệ số tự do, đặt nhân tử chung $3x$: $3x(x - 4) = 0$.

a) Một tích bằng $0$ khi một trong các thừa số bằng $0$: $3x = 0$ hoặc $x - 4 = 0$, suy ra $x = 0$ hoặc $x = 4$.

b) Phương trình khuyết hệ số của $x$; chuyển vế: $2x^2 = 18$, chia hai vế cho $2$: $x^2 = 9$.

b) Vì $9 > 0$ nên $x = 3$ hoặc $x = -3$ (hai số có bình phương bằng $9$).

c) Chuyển vế: $5x^2 = -20$, chia hai vế cho $5$: $x^2 = -4$.

c) Vì $x^2 \geq 0$ với mọi $x$ còn $-4 < 0$, nên không có giá trị nào của $x$ thoả mãn; phương trình vô nghiệm.

Đáp số: a) $\boxed{x = 0;\ x = 4}$; b) $\boxed{x = 3;\ x = -3}$; c) $\boxed{\text{Phương trình vô nghiệm}}$
\end{solution}
\end{problem}

% ---------------------------------------------------------------- công thức nghiệm, đưa về dạng tổng quát
\begin{problem}[9DS-PT2-002]
% nguon: dummy (tự sinh để test)
% lop: 9
% kc_du_kien: 9DS-PT2-02
% do_kho_du_kien: 2
Giải phương trình $x(2x - 3) = 2(x + 1) + 1$.

\begin{solution}
Khai triển hai vế (nhân đơn thức với đa thức): $2x^2 - 3x = 2x + 3$.

Chuyển tất cả các hạng tử sang vế trái để đưa về dạng $ax^2 + bx + c = 0$: $2x^2 - 5x - 3 = 0$, với $a = 2$, $b = -5$, $c = -3$.

Tính biệt thức: $\Delta = b^2 - 4ac = (-5)^2 - 4 \cdot 2 \cdot (-3) = 25 + 24 = 49 > 0$, nên phương trình có hai nghiệm phân biệt; $\sqrt{\Delta} = 7$.

Áp dụng công thức nghiệm: $x_1 = \dfrac{-b + \sqrt{\Delta}}{2a} = \dfrac{5 + 7}{4} = 3$ và $x_2 = \dfrac{-b - \sqrt{\Delta}}{2a} = \dfrac{5 - 7}{4} = -\dfrac{1}{2}$.

Thử lại: với $x = 3$, vế trái $3 \cdot 3 = 9$, vế phải $2 \cdot 4 + 1 = 9$ (đúng); với $x = -\dfrac{1}{2}$, vế trái $-\dfrac{1}{2} \cdot (-4) = 2$, vế phải $2 \cdot \dfrac{1}{2} + 1 = 2$ (đúng).

Đáp số: $\boxed{x = 3;\ x = -\dfrac{1}{2}}$
\end{solution}
\end{problem}

% ---------------------------------------------------------------- ba trường hợp của biệt thức
\begin{problem}[9DS-PT2-003]
% nguon: dummy (tự sinh để test)
% lop: 9
% kc_du_kien: 9DS-PT2-02
% do_kho_du_kien: 2
Dùng công thức nghiệm (hoặc công thức nghiệm thu gọn), giải các phương trình sau:

a) $x^2 - 6x + 4 = 0$;

b) $9x^2 - 12x + 4 = 0$;

c) $2x^2 + 3x + 5 = 0$.

\begin{solution}
a) Hệ số $b = -6$ là số chẵn nên dùng công thức nghiệm thu gọn với $b' = -3$: $\Delta' = b'^2 - ac = (-3)^2 - 1 \cdot 4 = 5 > 0$, phương trình có hai nghiệm phân biệt.

a) Theo công thức nghiệm thu gọn: $x_{1,2} = \dfrac{-b' \pm \sqrt{\Delta'}}{a} = 3 \pm \sqrt{5}$.

b) Hệ số $b = -12$ chẵn, $b' = -6$: $\Delta' = (-6)^2 - 9 \cdot 4 = 36 - 36 = 0$, phương trình có nghiệm kép.

b) Nghiệm kép: $x_1 = x_2 = -\dfrac{b'}{a} = \dfrac{6}{9} = \dfrac{2}{3}$ (phù hợp với $9x^2 - 12x + 4 = (3x - 2)^2$ theo hằng đẳng thức bình phương của một hiệu).

c) Với $a = 2$, $b = 3$, $c = 5$: $\Delta = b^2 - 4ac = 9 - 40 = -31 < 0$.

c) Vì $\Delta < 0$ nên phương trình vô nghiệm.

Đáp số: a) $\boxed{x = 3 + \sqrt{5};\ x = 3 - \sqrt{5}}$; b) $\boxed{x_1 = x_2 = \dfrac{2}{3}}$; c) $\boxed{\text{Phương trình vô nghiệm}}$
\end{solution}
\end{problem}

% ---------------------------------------------------------------- nhẩm nghiệm
\begin{problem}[9DS-PT2-004]
% nguon: dummy (tự sinh để test)
% lop: 9
% kc_du_kien: 9DS-PT2-04
% do_kho_du_kien: 2
Không dùng công thức nghiệm, hãy nhẩm nghiệm của các phương trình sau:

a) $5x^2 - 7x + 2 = 0$;

b) $2026x^2 + 2025x - 1 = 0$.

\begin{solution}
a) Phương trình có $a = 5$, $b = -7$, $c = 2$ và $a + b + c = 5 - 7 + 2 = 0$.

a) Vì $a + b + c = 0$ nên phương trình có hai nghiệm $x_1 = 1$ và $x_2 = \dfrac{c}{a} = \dfrac{2}{5}$.

b) Phương trình có $a = 2026$, $b = 2025$, $c = -1$ và $a - b + c = 2026 - 2025 - 1 = 0$.

b) Vì $a - b + c = 0$ nên phương trình có hai nghiệm $x_1 = -1$ và $x_2 = -\dfrac{c}{a} = \dfrac{1}{2026}$.

Đáp số: a) $\boxed{x_1 = 1;\ x_2 = \dfrac{2}{5}}$; b) $\boxed{x_1 = -1;\ x_2 = \dfrac{1}{2026}}$
\end{solution}
\end{problem}

% ---------------------------------------------------------------- Viète: biểu thức đối xứng
\begin{problem}[9DS-PT2-005]
% nguon: dummy (tự sinh để test)
% lop: 9
% kc_du_kien: 9DS-PT2-03, 9DS-PT2-02, 8DS-HDT-04
% do_kho_du_kien: 3
Cho phương trình $x^2 - 5x + 3 = 0$.

a) Chứng tỏ rằng phương trình có hai nghiệm phân biệt $x_1, x_2$.

b) Không giải phương trình, hãy tính giá trị của các biểu thức $A = x_1^2 + x_2^2$, $B = \dfrac{1}{x_1} + \dfrac{1}{x_2}$ và $C = x_1^3 + x_2^3$.

\begin{solution}
a) Phương trình có $a = 1$, $b = -5$, $c = 3$: $\Delta = (-5)^2 - 4 \cdot 1 \cdot 3 = 25 - 12 = 13 > 0$, nên phương trình có hai nghiệm phân biệt $x_1, x_2$.

b) Theo định lý Viète: $x_1 + x_2 = -\dfrac{b}{a} = 5$ và $x_1 x_2 = \dfrac{c}{a} = 3$.

b) Thêm bớt $2x_1x_2$ để dùng hằng đẳng thức bình phương của một tổng: $A = (x_1 + x_2)^2 - 2x_1x_2 = 5^2 - 2 \cdot 3 = 19$.

b) Vì $x_1 x_2 = 3 \neq 0$ nên $x_1 \neq 0$, $x_2 \neq 0$ và $B$ xác định; quy đồng mẫu: $B = \dfrac{x_1 + x_2}{x_1 x_2} = \dfrac{5}{3}$.

b) Dùng hằng đẳng thức tổng hai lập phương: $C = (x_1 + x_2)(x_1^2 - x_1x_2 + x_2^2) = (x_1 + x_2)(A - x_1x_2) = 5 \cdot (19 - 3) = 80$.

Đáp số: a) $\boxed{\Delta = 13 > 0}$; b) $\boxed{A = 19,\ B = \dfrac{5}{3},\ C = 80}$
\end{solution}
\end{problem}

% ---------------------------------------------------------------- tìm hai số biết tổng và tích
\begin{problem}[9DS-PT2-006]
% nguon: dummy (tự sinh để test)
% lop: 9
% kc_du_kien: 9DS-PT2-05, 9DS-PT2-02
% do_kho_du_kien: 3
a) Tìm hai số biết tổng của chúng bằng $17$ và tích của chúng bằng $72$.

b) Tìm hai số $u$, $v$ biết $u - v = 3$ và $uv = 40$.

c) Có tồn tại hai số có tổng bằng $3$ và tích bằng $5$ hay không? Vì sao?

\begin{solution}
a) Với $S = 17$, $P = 72$ ta có $S^2 - 4P = 289 - 288 = 1 \geq 0$, nên hai số cần tìm tồn tại và là hai nghiệm của phương trình $X^2 - 17X + 72 = 0$.

a) Giải phương trình: $\Delta = 17^2 - 4 \cdot 72 = 1$, $\sqrt{\Delta} = 1$, nên $X_1 = \dfrac{17 + 1}{2} = 9$ và $X_2 = \dfrac{17 - 1}{2} = 8$.

a) Thử lại: $9 + 8 = 17$ và $9 \cdot 8 = 72$ (đúng). Vậy hai số cần tìm là $8$ và $9$.

b) Đặt $t = -v$. Khi đó $u + t = u - v = 3$ và $ut = -uv = -40$, nên $u$ và $t$ là hai nghiệm của phương trình $X^2 - 3X - 40 = 0$ (điều kiện $S^2 - 4P = 9 + 160 = 169 \geq 0$ thoả mãn).

b) Giải phương trình: $\Delta = 9 + 160 = 169$, $\sqrt{\Delta} = 13$, nên $X_1 = \dfrac{3 + 13}{2} = 8$ và $X_2 = \dfrac{3 - 13}{2} = -5$.

b) Nếu $u = 8$, $t = -5$ thì $v = 5$; nếu $u = -5$, $t = 8$ thì $v = -8$. Thử lại: $8 - 5 = 3$, $8 \cdot 5 = 40$ và $-5 - (-8) = 3$, $(-5) \cdot (-8) = 40$ (đúng).

c) Nếu có hai số như vậy thì chúng là hai nghiệm của phương trình $X^2 - 3X + 5 = 0$.

c) Phương trình này có $\Delta = 9 - 20 = -11 < 0$ (tức là $S^2 < 4P$) nên vô nghiệm; do đó không tồn tại hai số như vậy.

Đáp số: a) $\boxed{8 \text{ và } 9}$; b) $\boxed{(u; v) = (8; 5) \text{ hoặc } (u; v) = (-5; -8)}$; c) $\boxed{\text{Không tồn tại}}$
\end{solution}
\end{problem}

% ---------------------------------------------------------------- tham số: hai nghiệm phân biệt, hệ thức, loại m
\begin{problem}[9DS-PT2-007]
% nguon: dummy (tự sinh để test)
% lop: 9
% kc_du_kien: 9DS-PT2-06, 9DS-PT2-03
% do_kho_du_kien: 4
Cho phương trình $x^2 - 2(m + 1)x + m^2 + 3 = 0$ $(1)$, với $m$ là tham số.

a) Tìm $m$ để phương trình $(1)$ có hai nghiệm phân biệt.

b) Tìm $m$ để phương trình $(1)$ có hai nghiệm $x_1, x_2$ thoả mãn $x_1^2 + x_2^2 = 22$.

\begin{solution}
a) Phương trình $(1)$ có $a = 1$, $b' = -(m + 1)$, $c = m^2 + 3$, nên $\Delta' = (m + 1)^2 - (m^2 + 3) = m^2 + 2m + 1 - m^2 - 3 = 2m - 2$.

a) Phương trình $(1)$ có hai nghiệm phân biệt khi và chỉ khi $\Delta' > 0 \Leftrightarrow 2m - 2 > 0 \Leftrightarrow m > 1$.

b) Phương trình $(1)$ có hai nghiệm (phân biệt hoặc kép) khi và chỉ khi $\Delta' \geq 0 \Leftrightarrow m \geq 1$. Khi đó, theo định lý Viète: $x_1 + x_2 = 2(m + 1)$ và $x_1 x_2 = m^2 + 3$.

b) Biến đổi theo tổng và tích: $x_1^2 + x_2^2 = (x_1 + x_2)^2 - 2x_1x_2 = 4(m + 1)^2 - 2(m^2 + 3) = 2m^2 + 8m - 2$.

b) Điều kiện đề bài: $2m^2 + 8m - 2 = 22 \Leftrightarrow m^2 + 4m - 12 = 0$; phương trình này có $\Delta' = 2^2 + 12 = 16$, nên $m = -2 + 4 = 2$ hoặc $m = -2 - 4 = -6$.

b) Đối chiếu điều kiện $m \geq 1$: $m = 2$ thoả mãn; $m = -6$ bị loại (khi đó $\Delta' = -14 < 0$, phương trình $(1)$ vô nghiệm).

b) Thử lại với $m = 2$: phương trình $x^2 - 6x + 7 = 0$ có $\Delta' = 2 > 0$, $x_1 + x_2 = 6$, $x_1x_2 = 7$, nên $x_1^2 + x_2^2 = 36 - 14 = 22$ (đúng).

Đáp số: a) $\boxed{m > 1}$; b) $\boxed{m = 2}$
\end{solution}
\end{problem}

% ---------------------------------------------------------------- tham số, hệ thức không đối xứng, hai cách
\begin{problem}[9DS-PT2-008]
% nguon: dummy (tự sinh để test)
% lop: 9
% kc_du_kien: 9DS-PT2-06, 9DS-PT2-03, 9DS-PT2-04
% do_kho_du_kien: 4
Cho phương trình $x^2 - mx + m - 1 = 0$ $(1)$, với $m$ là tham số.

a) Chứng minh rằng phương trình $(1)$ luôn có nghiệm với mọi giá trị của $m$.

b) Tìm $m$ để phương trình $(1)$ có hai nghiệm $x_1, x_2$ thoả mãn $x_1 = 2x_2$.

\begin{solution}[cách 1: dùng biệt thức và định lý Viète]
a) Phương trình $(1)$ có $a = 1$, $b = -m$, $c = m - 1$, nên $\Delta = m^2 - 4(m - 1) = m^2 - 4m + 4 = (m - 2)^2$ (hằng đẳng thức bình phương của một hiệu).

a) Vì $(m - 2)^2 \geq 0$ với mọi $m$ nên $\Delta \geq 0$; do đó phương trình $(1)$ luôn có nghiệm với mọi $m$.

b) Theo câu a), phương trình $(1)$ luôn có hai nghiệm $x_1, x_2$ (có thể trùng nhau). Theo định lý Viète: $x_1 + x_2 = m$ và $x_1 x_2 = m - 1$.

b) Kết hợp $x_1 = 2x_2$ với $x_1 + x_2 = m$: $3x_2 = m$, suy ra $x_2 = \dfrac{m}{3}$ và $x_1 = \dfrac{2m}{3}$.

b) Thay vào $x_1 x_2 = m - 1$: $\dfrac{2m^2}{9} = m - 1 \Leftrightarrow 2m^2 - 9m + 9 = 0$; phương trình này có $\Delta = 81 - 72 = 9$, nên $m = \dfrac{9 + 3}{4} = 3$ hoặc $m = \dfrac{9 - 3}{4} = \dfrac{3}{2}$.

b) Thử lại: với $m = 3$, phương trình $x^2 - 3x + 2 = 0$ có hai nghiệm $2$ và $1$, và $2 = 2 \cdot 1$ (đúng); với $m = \dfrac{3}{2}$, phương trình tương đương $2x^2 - 3x + 1 = 0$ có hai nghiệm $1$ và $\dfrac{1}{2}$, và $1 = 2 \cdot \dfrac{1}{2}$ (đúng).

Đáp số: a) $\boxed{\Delta = (m - 2)^2 \geq 0 \text{ với mọi } m}$; b) $\boxed{m = 3 \text{ hoặc } m = \dfrac{3}{2}}$
\end{solution}

\begin{solution}[cách 2: nhẩm nghiệm]
a) Phương trình $(1)$ có $a + b + c = 1 + (-m) + (m - 1) = 0$ với mọi $m$.

a) Do đó $(1)$ luôn có hai nghiệm $x = 1$ và $x = \dfrac{c}{a} = m - 1$, tức là luôn có nghiệm với mọi $m$.

b) Hai nghiệm của $(1)$ là $1$ và $m - 1$; vì vai trò của $x_1, x_2$ trong hệ thức không như nhau nên xét hai trường hợp.

b) Trường hợp $x_1 = 1$, $x_2 = m - 1$: $1 = 2(m - 1) \Leftrightarrow m = \dfrac{3}{2}$.

b) Trường hợp $x_1 = m - 1$, $x_2 = 1$: $m - 1 = 2 \Leftrightarrow m = 3$.

Đáp số: a) $\boxed{\text{Phương trình } (1) \text{ luôn có nghiệm } x = 1}$; b) $\boxed{m = 3 \text{ hoặc } m = \dfrac{3}{2}}$
\end{solution}
\end{problem}

% ---------------------------------------------------------------- tham số: thay nghiệm, giá trị nhỏ nhất
\begin{problem}[9DS-PT2-009]
% nguon: dummy (tự sinh để test)
% lop: 9
% kc_du_kien: 9DS-PT2-06, 9DS-PT2-03
% do_kho_du_kien: 5
Cho phương trình $x^2 - 2(m + 1)x + 2m = 0$ $(1)$, với $m$ là tham số.

a) Chứng minh rằng phương trình $(1)$ luôn có hai nghiệm phân biệt với mọi $m$.

b) Gọi $x_1, x_2$ là hai nghiệm của $(1)$. Tìm $m$ để $x_1^2 + 2(m + 1)x_2 = 8$.

c) Tìm giá trị nhỏ nhất của biểu thức $A = x_1^2 + x_2^2$ và giá trị của $m$ tương ứng.

\begin{solution}
a) Phương trình $(1)$ có $a = 1$, $b' = -(m + 1)$, $c = 2m$, nên $\Delta' = (m + 1)^2 - 2m = m^2 + 1$.

a) Vì $m^2 \geq 0$ nên $\Delta' = m^2 + 1 \geq 1 > 0$ với mọi $m$; vậy phương trình $(1)$ luôn có hai nghiệm phân biệt.

b) Theo định lý Viète: $x_1 + x_2 = 2(m + 1)$ và $x_1 x_2 = 2m$.

b) Vì $x_1$ là nghiệm của $(1)$ nên $x_1^2 - 2(m + 1)x_1 + 2m = 0$, suy ra $x_1^2 = 2(m + 1)x_1 - 2m$.

b) Thay vào biểu thức: $x_1^2 + 2(m + 1)x_2 = 2(m + 1)x_1 - 2m + 2(m + 1)x_2 = 2(m + 1)(x_1 + x_2) - 2m = 2(m + 1) \cdot 2(m + 1) - 2m = 4m^2 + 6m + 4$.

b) Điều kiện đề bài: $4m^2 + 6m + 4 = 8 \Leftrightarrow 2m^2 + 3m - 2 = 0$; phương trình này có $\Delta = 9 + 16 = 25$, nên $m = \dfrac{-3 + 5}{4} = \dfrac{1}{2}$ hoặc $m = \dfrac{-3 - 5}{4} = -2$.

b) Cả hai giá trị đều thoả mãn vì theo câu a) phương trình $(1)$ có hai nghiệm phân biệt với mọi $m$. Thử lại với $m = \dfrac{1}{2}$: $(1)$ là $x^2 - 3x + 1 = 0$, có $x_1^2 = 3x_1 - 1$ nên $x_1^2 + 3x_2 = 3(x_1 + x_2) - 1 = 9 - 1 = 8$ (đúng).

c) Biến đổi theo tổng và tích: $A = (x_1 + x_2)^2 - 2x_1x_2 = 4(m + 1)^2 - 4m = 4m^2 + 4m + 4$.

c) Tách thành bình phương của một tổng: $A = (2m + 1)^2 + 3 \geq 3$ với mọi $m$, vì $(2m + 1)^2 \geq 0$.

c) Dấu bằng xảy ra khi và chỉ khi $2m + 1 = 0 \Leftrightarrow m = -\dfrac{1}{2}$. Vậy giá trị nhỏ nhất của $A$ là $3$, đạt được khi $m = -\dfrac{1}{2}$.

Đáp số: a) $\boxed{\Delta' = m^2 + 1 > 0 \text{ với mọi } m}$; b) $\boxed{m = \dfrac{1}{2} \text{ hoặc } m = -2}$; c) $\boxed{A_{\min} = 3 \text{ khi } m = -\dfrac{1}{2}}$
\end{solution}
\end{problem}

% ---------------------------------------------------------------- bài toán chuyển động
\begin{problem}[9DS-PT2-010]
% nguon: dummy (tự sinh để test)
% lop: 9
% kc_du_kien: 9DS-PT2-07, 9DS-PT2-02
% do_kho_du_kien: 3
Một ô tô đi từ $A$ đến $B$ cách nhau $120$ km với vận tốc không đổi. Lúc về, ô tô tăng vận tốc thêm $10$ km/h so với lúc đi nên thời gian về ít hơn thời gian đi $1$ giờ. Tính vận tốc của ô tô lúc đi.

\begin{solution}
Gọi vận tốc của ô tô lúc đi là $x$ (km/h), điều kiện $x > 0$.

Thời gian đi là $\dfrac{120}{x}$ (giờ). Vận tốc lúc về là $x + 10$ (km/h), nên thời gian về là $\dfrac{120}{x + 10}$ (giờ).

Thời gian về ít hơn thời gian đi $1$ giờ nên ta có phương trình $\dfrac{120}{x} - \dfrac{120}{x + 10} = 1$.

Vì $x > 0$ nên $x(x + 10) \neq 0$; quy đồng mẫu và khử mẫu: $120(x + 10) - 120x = x(x + 10) \Leftrightarrow 1200 = x^2 + 10x \Leftrightarrow x^2 + 10x - 1200 = 0$.

Giải bằng công thức nghiệm thu gọn với $b' = 5$: $\Delta' = 5^2 + 1200 = 1225$, $\sqrt{\Delta'} = 35$, nên $x_1 = -5 + 35 = 30$ và $x_2 = -5 - 35 = -40$.

Đối chiếu điều kiện $x > 0$: $x_1 = 30$ thoả mãn, $x_2 = -40$ bị loại (vận tốc phải dương).

Thử lại: thời gian đi $\dfrac{120}{30} = 4$ giờ, thời gian về $\dfrac{120}{40} = 3$ giờ, ít hơn $1$ giờ (đúng).

Đáp số: $\boxed{30 \text{ km/h}}$
\end{solution}
\end{problem}

% ---------------------------------------------------------------- biến đổi về phương trình bậc hai khuyết
\begin{problem}[9DS-PT2-011]
% nguon: dummy (tự sinh để test)
% lop: 9
% kc_du_kien: 9DS-PT2-01, 8DS-HDT-01, 8DS-HDT-02
% do_kho_du_kien: 1
Đưa các phương trình sau về dạng $ax^2 + bx + c = 0$ rồi giải:

a) $(x - 1)(x + 1) = 2x^2 - 10$;

b) $(x + 3)^2 = 9 - 2x$.

\begin{solution}
a) Theo hằng đẳng thức hiệu hai bình phương, vế trái bằng $x^2 - 1$; phương trình trở thành $x^2 - 1 = 2x^2 - 10$. Chuyển vế và thu gọn: $-x^2 + 9 = 0$, tức là $x^2 - 9 = 0$ (phương trình bậc hai khuyết hệ số $b$).

a) Chuyển vế: $x^2 = 9$; vì $9 > 0$ nên $x = 3$ hoặc $x = -3$ (hai số có bình phương bằng $9$).

b) Khai triển vế trái theo hằng đẳng thức bình phương của một tổng: $x^2 + 6x + 9 = 9 - 2x$. Chuyển vế và thu gọn: $x^2 + 8x = 0$ (phương trình bậc hai khuyết hệ số tự do).

b) Đặt nhân tử chung (không chia hai vế cho $x$ để tránh mất nghiệm $x = 0$): $x(x + 8) = 0$, nên $x = 0$ hoặc $x + 8 = 0$, tức là $x = 0$ hoặc $x = -8$.

Đáp số: a) $\boxed{x = 3;\ x = -3}$; b) $\boxed{x = 0;\ x = -8}$
\end{solution}
\end{problem}

% ---------------------------------------------------------------- bài toán thực tế: lãi kép
\begin{problem}[9DS-PT2-012]
% nguon: dummy (tự sinh để test)
% lop: 9
% kc_du_kien: 9DS-PT2-07, 9DS-PT2-01
% do_kho_du_kien: 2
Bác Lan gửi $100\,000\,000$ đồng vào ngân hàng theo thể thức lãi kép (tiền lãi của mỗi năm được nhập vào vốn để tính lãi cho năm sau), lãi suất mỗi năm không đổi. Sau đúng $2$ năm, bác rút được cả vốn lẫn lãi là $112\,360\,000$ đồng. Tính lãi suất mỗi năm của ngân hàng.

\begin{solution}
Gọi lãi suất mỗi năm là $x$ (viết dưới dạng phân số, chẳng hạn $x = \dfrac{1}{10}$ ứng với $10\%$), điều kiện $x > 0$.

Sau năm thứ nhất, số tiền cả vốn lẫn lãi là $100\,000\,000 + 100\,000\,000x = 100\,000\,000(1 + x)$ (đồng). Số tiền này là vốn của năm thứ hai, nên sau năm thứ hai số tiền là $100\,000\,000(1 + x)(1 + x) = 100\,000\,000(1 + x)^2$ (đồng).

Theo đề bài ta có phương trình $100\,000\,000(1 + x)^2 = 112\,360\,000$. Chia hai vế cho $100\,000\,000$: $(1 + x)^2 = \dfrac{11236}{10000}$.

Vì $106^2 = 11236$ và $100^2 = 10000$ nên $(1 + x)^2 = \left(\dfrac{106}{100}\right)^2$. Hai số có bình phương bằng nhau thì bằng nhau hoặc đối nhau: $1 + x = \dfrac{106}{100}$ hoặc $1 + x = -\dfrac{106}{100}$.

Suy ra $x = \dfrac{6}{100}$ hoặc $x = -\dfrac{206}{100}$. Đối chiếu điều kiện $x > 0$: $x = \dfrac{6}{100} = 6\%$ thoả mãn; $x = -\dfrac{206}{100}$ bị loại.

Thử lại: sau năm thứ nhất bác có $100\,000\,000 \cdot \dfrac{106}{100} = 106\,000\,000$ đồng; sau năm thứ hai có $106\,000\,000 \cdot \dfrac{106}{100} = 112\,360\,000$ đồng (đúng).

Đáp số: $\boxed{6\% \text{ mỗi năm}}$
\end{solution}
\end{problem}

% ---------------------------------------------------------------- biết một nghiệm, tìm tham số và nghiệm còn lại, hai cách
\begin{problem}[9DS-PT2-013]
% nguon: dummy (tự sinh để test)
% lop: 9
% kc_du_kien: 9DS-PT2-06, 9DS-PT2-03, 9DS-PT2-02
% do_kho_du_kien: 2
Cho phương trình $x^2 - (m + 2)x + 2m + 4 = 0$ $(1)$, với $m$ là tham số. Biết phương trình $(1)$ có một nghiệm là $x = 3$. Tìm $m$ và nghiệm còn lại của phương trình.

\begin{solution}[cách 1: dùng định lý Viète]
Vì $x = 3$ là nghiệm của $(1)$ nên khi thay $x = 3$ vào $(1)$ ta được đẳng thức đúng: $9 - 3(m + 2) + 2m + 4 = 0 \Leftrightarrow 7 - m = 0 \Leftrightarrow m = 7$.

Với $m = 7$, phương trình $(1)$ có nghiệm $x_1 = 3$ nên có hai nghiệm $x_1, x_2$ (có thể trùng nhau); theo định lý Viète: $x_1 + x_2 = m + 2 = 9$.

Suy ra nghiệm còn lại $x_2 = 9 - 3 = 6$. Kiểm tra bằng tích: $x_1 x_2 = 3 \cdot 6 = 18 = 2 \cdot 7 + 4$ (đúng).

Đáp số: $\boxed{m = 7;\ \text{nghiệm còn lại } x = 6}$
\end{solution}

\begin{solution}[cách 2: thay $m$ rồi giải phương trình]
Thay $x = 3$ vào $(1)$: $9 - 3(m + 2) + 2m + 4 = 0 \Leftrightarrow 7 - m = 0 \Leftrightarrow m = 7$.

Với $m = 7$, phương trình $(1)$ trở thành $x^2 - 9x + 18 = 0$, có $\Delta = (-9)^2 - 4 \cdot 1 \cdot 18 = 81 - 72 = 9 > 0$, $\sqrt{\Delta} = 3$.

Theo công thức nghiệm: $x_1 = \dfrac{9 + 3}{2} = 6$ và $x_2 = \dfrac{9 - 3}{2} = 3$; nghiệm còn lại (khác $3$) là $x = 6$.

Đáp số: $\boxed{m = 7;\ \text{nghiệm còn lại } x = 6}$
\end{solution}
\end{problem}

% ---------------------------------------------------------------- phương trình trùng phương
\begin{problem}[9DS-PT2-014]
% nguon: dummy (tự sinh để test)
% lop: 9
% kc_du_kien: 9DS-PT2-08, 9DS-PT2-02, 9DS-PT2-01
% do_kho_du_kien: 3
Giải phương trình $x^4 - 5x^2 - 36 = 0$.

\begin{solution}
Đặt $t = x^2$, điều kiện $t \geq 0$ (bình phương của một số luôn không âm). Phương trình trở thành phương trình bậc hai ẩn $t$: $t^2 - 5t - 36 = 0$.

Với $a = 1$, $b = -5$, $c = -36$: $\Delta = (-5)^2 - 4 \cdot 1 \cdot (-36) = 25 + 144 = 169 > 0$, $\sqrt{\Delta} = 13$.

Theo công thức nghiệm: $t_1 = \dfrac{5 + 13}{2} = 9$ và $t_2 = \dfrac{5 - 13}{2} = -4$.

Đối chiếu điều kiện $t \geq 0$: $t_1 = 9$ thoả mãn; $t_2 = -4$ bị loại (không có $x$ nào để $x^2 = -4$).

Với $t = 9$: $x^2 = 9$, suy ra $x = 3$ hoặc $x = -3$.

Thử lại: $3^4 - 5 \cdot 3^2 - 36 = 81 - 45 - 36 = 0$; vì phương trình chỉ chứa luỹ thừa bậc chẵn của $x$ nên $x = -3$ cũng thoả mãn.

Đáp số: $\boxed{x = 3;\ x = -3}$
\end{solution}
\end{problem}

% ---------------------------------------------------------------- phương trình chứa ẩn ở mẫu dẫn tới bậc hai
\begin{problem}[9DS-PT2-015]
% nguon: dummy (tự sinh để test)
% lop: 9
% kc_du_kien: 9DS-PT2-08, 9DS-PT2-04, 8DS-PS-03
% do_kho_du_kien: 3
Giải phương trình $\dfrac{x}{x - 1} + \dfrac{2}{x + 1} = \dfrac{2}{x^2 - 1}$.

\begin{solution}
Điều kiện xác định: các mẫu khác $0$, tức là $x - 1 \neq 0$, $x + 1 \neq 0$ và $x^2 - 1 = (x - 1)(x + 1) \neq 0$; do đó $x \neq 1$ và $x \neq -1$.

Theo hằng đẳng thức hiệu hai bình phương, mẫu chung là $(x - 1)(x + 1) = x^2 - 1$. Quy đồng mẫu hai vế rồi khử mẫu: $x(x + 1) + 2(x - 1) = 2$.

Khai triển và thu gọn: $x^2 + x + 2x - 2 - 2 = 0$, tức là $x^2 + 3x - 4 = 0$.

Phương trình này có $a = 1$, $b = 3$, $c = -4$ và $a + b + c = 1 + 3 - 4 = 0$, nên có hai nghiệm $x_1 = 1$ và $x_2 = \dfrac{c}{a} = -4$.

Đối chiếu điều kiện: $x_1 = 1$ không thoả $x \neq 1$ nên loại; $x_2 = -4$ thoả mãn.

Thử lại với $x = -4$: vế trái $\dfrac{-4}{-5} + \dfrac{2}{-3} = \dfrac{4}{5} - \dfrac{2}{3} = \dfrac{2}{15}$, vế phải $\dfrac{2}{16 - 1} = \dfrac{2}{15}$ (đúng).

Đáp số: $\boxed{x = -4}$
\end{solution}
\end{problem}

% ---------------------------------------------------------------- nghiệm trái dấu, lập phương trình bậc hai có nghiệm cho trước
\begin{problem}[9DS-PT2-016]
% nguon: dummy (tự sinh để test)
% lop: 9
% kc_du_kien: 9DS-PT2-05, 9DS-PT2-03
% do_kho_du_kien: 3
Cho phương trình $x^2 - 3x - 5 = 0$.

a) Không giải phương trình, chứng tỏ phương trình có hai nghiệm trái dấu $x_1, x_2$.

b) Không giải phương trình, hãy lập một phương trình bậc hai ẩn $y$ có hai nghiệm là $y_1 = x_1 + \dfrac{1}{x_2}$ và $y_2 = x_2 + \dfrac{1}{x_1}$.

\begin{solution}
a) Phương trình có $a = 1$, $b = -3$, $c = -5$ nên $ac = -5 < 0$. Khi đó $\Delta = b^2 - 4ac > 0$ (vì $b^2 \geq 0$ và $-4ac > 0$), nên phương trình có hai nghiệm phân biệt $x_1, x_2$.

a) Theo định lý Viète, $x_1 x_2 = \dfrac{c}{a} = -5 < 0$, nên hai nghiệm trái dấu.

b) Theo định lý Viète: $x_1 + x_2 = 3$ và $x_1 x_2 = -5$; vì $x_1 x_2 \neq 0$ nên $x_1 \neq 0$, $x_2 \neq 0$ và $y_1, y_2$ xác định.

b) Tính tổng, quy đồng $\dfrac{1}{x_1} + \dfrac{1}{x_2}$: $y_1 + y_2 = (x_1 + x_2) + \dfrac{x_1 + x_2}{x_1 x_2} = 3 + \dfrac{3}{-5} = \dfrac{12}{5}$.

b) Tính tích bằng cách nhân phá ngoặc: $y_1 y_2 = \left(x_1 + \dfrac{1}{x_2}\right)\left(x_2 + \dfrac{1}{x_1}\right) = x_1x_2 + 1 + 1 + \dfrac{1}{x_1x_2} = -5 + 2 - \dfrac{1}{5} = -\dfrac{16}{5}$.

b) Hai số $y_1, y_2$ có tổng $S = \dfrac{12}{5}$ và tích $P = -\dfrac{16}{5}$ nên là hai nghiệm của phương trình $y^2 - Sy + P = 0$, tức là $y^2 - \dfrac{12}{5}y - \dfrac{16}{5} = 0$; nhân hai vế với $5$ được $5y^2 - 12y - 16 = 0$.

Đáp số: a) $\boxed{x_1 x_2 = -5 < 0}$; b) $\boxed{5y^2 - 12y - 16 = 0}$
\end{solution}
\end{problem}

% ---------------------------------------------------------------- tham số: hai nghiệm trái dấu, so sánh giá trị tuyệt đối
\begin{problem}[9DS-PT2-017]
% nguon: dummy (tự sinh để test)
% lop: 9
% kc_du_kien: 9DS-PT2-06, 9DS-PT2-03, 9DS-BPT-02
% do_kho_du_kien: 4
Cho phương trình $x^2 - 2(m - 3)x + m - 5 = 0$ $(1)$, với $m$ là tham số.

a) Giải phương trình $(1)$ khi $m = 4$.

b) Tìm $m$ để phương trình $(1)$ có hai nghiệm trái dấu, trong đó nghiệm âm có giá trị tuyệt đối lớn hơn nghiệm dương.

\begin{solution}
a) Với $m = 4$, phương trình $(1)$ trở thành $x^2 - 2x - 1 = 0$, có $a = 1$, $b' = -1$, $c = -1$ và $\Delta' = (-1)^2 - 1 \cdot (-1) = 2 > 0$.

a) Theo công thức nghiệm thu gọn: $x = \dfrac{-b' \pm \sqrt{\Delta'}}{a} = 1 \pm \sqrt{2}$.

b) Phương trình $(1)$ có $a = 1$, $c = m - 5$. Nếu $ac < 0$ thì $\Delta' = b'^2 - ac > 0$ và theo định lý Viète $x_1x_2 = \dfrac{c}{a} < 0$; ngược lại, nếu hai nghiệm trái dấu thì $\dfrac{c}{a} = x_1x_2 < 0$. Do đó $(1)$ có hai nghiệm trái dấu khi và chỉ khi $ac < 0 \Leftrightarrow m - 5 < 0 \Leftrightarrow m < 5$.

b) Khi đó gọi $x_1 < 0 < x_2$. Nghiệm âm có giá trị tuyệt đối lớn hơn nghiệm dương: $|x_1| > x_2 \Leftrightarrow -x_1 > x_2 \Leftrightarrow x_1 + x_2 < 0$.

b) Theo định lý Viète: $x_1 + x_2 = 2(m - 3)$, nên điều kiện trở thành $2(m - 3) < 0 \Leftrightarrow m < 3$.

b) Kết hợp với điều kiện $m < 5$ ta được $m < 3$.

Đáp số: a) $\boxed{x = 1 + \sqrt{2};\ x = 1 - \sqrt{2}}$; b) $\boxed{m < 3}$
\end{solution}
\end{problem}

% ---------------------------------------------------------------- tham số: hệ thức không phụ thuộc m, hiệu hai nghiệm
\begin{problem}[9DS-PT2-018]
% nguon: dummy (tự sinh để test)
% lop: 9
% kc_du_kien: 9DS-PT2-06, 9DS-PT2-03, 8DS-HDT-01
% do_kho_du_kien: 4
Cho phương trình $x^2 - (m + 3)x + 2m + 2 = 0$ $(1)$, với $m$ là tham số.

a) Chứng minh rằng phương trình $(1)$ luôn có nghiệm với mọi $m$.

b) Gọi $x_1, x_2$ là hai nghiệm của $(1)$. Tìm một hệ thức liên hệ giữa $x_1$ và $x_2$ không phụ thuộc vào $m$.

c) Tìm $m$ để $|x_1 - x_2| = 3$.

\begin{solution}
a) Phương trình $(1)$ có $a = 1$, $b = -(m + 3)$, $c = 2m + 2$, nên $\Delta = (m + 3)^2 - 4(2m + 2) = m^2 + 6m + 9 - 8m - 8 = m^2 - 2m + 1 = (m - 1)^2$ (hằng đẳng thức bình phương của một hiệu).

a) Vì $(m - 1)^2 \geq 0$ với mọi $m$ nên $\Delta \geq 0$; phương trình $(1)$ luôn có nghiệm với mọi $m$.

b) Theo định lý Viète: $x_1 + x_2 = m + 3$ và $x_1 x_2 = 2m + 2$.

b) Từ hệ thức thứ nhất: $m = x_1 + x_2 - 3$. Thay vào hệ thức thứ hai để khử $m$: $x_1 x_2 = 2(x_1 + x_2 - 3) + 2 = 2(x_1 + x_2) - 4$, tức là $x_1 x_2 - 2(x_1 + x_2) + 4 = 0$.

c) Vì $|x_1 - x_2| \geq 0$ và $3 > 0$ nên $|x_1 - x_2| = 3 \Leftrightarrow (x_1 - x_2)^2 = 9$.

c) Biến đổi theo tổng và tích: $(x_1 - x_2)^2 = (x_1 + x_2)^2 - 4x_1x_2 = (m + 3)^2 - 4(2m + 2) = (m - 1)^2$ (theo câu a).

c) Do đó $(m - 1)^2 = 9 \Leftrightarrow m - 1 = 3$ hoặc $m - 1 = -3 \Leftrightarrow m = 4$ hoặc $m = -2$. Cả hai giá trị đều nhận vì theo câu a) phương trình $(1)$ luôn có nghiệm.

c) Thử lại: với $m = 4$, $(1)$ là $x^2 - 7x + 10 = 0$ có hai nghiệm $2$ và $5$, $|2 - 5| = 3$ (đúng); với $m = -2$, $(1)$ là $x^2 - x - 2 = 0$ có hai nghiệm $2$ và $-1$, $|2 - (-1)| = 3$ (đúng).

Đáp số: a) $\boxed{\Delta = (m - 1)^2 \geq 0 \text{ với mọi } m}$; b) $\boxed{x_1x_2 - 2(x_1 + x_2) + 4 = 0}$; c) $\boxed{m = 4 \text{ hoặc } m = -2}$
\end{solution}
\end{problem}

% ---------------------------------------------------------------- tích bốn thừa số, ghép cặp và đặt ẩn phụ
\begin{problem}[9DS-PT2-019]
% nguon: dummy (tự sinh để test)
% lop: 9
% kc_du_kien: 9DS-PT2-08, 9DS-PT2-02
% do_kho_du_kien: 4
Giải phương trình $(x - 1)(x + 1)(x + 3)(x + 5) = 105$.

\begin{solution}
Ghép thừa số thứ nhất với thứ tư, thứ hai với thứ ba (hai cặp có cùng tổng các số hạng tự do $-1 + 5 = 1 + 3$): $(x - 1)(x + 5) = x^2 + 4x - 5$ và $(x + 1)(x + 3) = x^2 + 4x + 3$. Phương trình trở thành $(x^2 + 4x - 5)(x^2 + 4x + 3) = 105$.

Đặt $t = x^2 + 4x$. Phương trình trở thành $(t - 5)(t + 3) = 105 \Leftrightarrow t^2 - 2t - 15 - 105 = 0 \Leftrightarrow t^2 - 2t - 120 = 0$.

Giải phương trình ẩn $t$ bằng công thức nghiệm thu gọn với $b' = -1$: $\Delta' = (-1)^2 - 1 \cdot (-120) = 121$, $\sqrt{\Delta'} = 11$, nên $t_1 = 1 + 11 = 12$ và $t_2 = 1 - 11 = -10$.

Với $t = 12$: $x^2 + 4x - 12 = 0$, có $\Delta' = 2^2 + 12 = 16$, $\sqrt{\Delta'} = 4$, nên $x = -2 + 4 = 2$ hoặc $x = -2 - 4 = -6$.

Với $t = -10$: $x^2 + 4x + 10 = 0$, có $\Delta' = 2^2 - 10 = -6 < 0$, nên phương trình này vô nghiệm.

Thử lại: với $x = 2$, vế trái $1 \cdot 3 \cdot 5 \cdot 7 = 105$; với $x = -6$, vế trái $(-7) \cdot (-5) \cdot (-3) \cdot (-1) = 105$ (đúng).

Đáp số: $\boxed{x = 2;\ x = -6}$
\end{solution}
\end{problem}

% ---------------------------------------------------------------- tham số: hai nghiệm là độ dài hai cạnh góc vuông
\begin{problem}[9DS-PT2-020]
% nguon: dummy (tự sinh để test)
% lop: 9
% kc_du_kien: 9DS-PT2-06, 9DS-PT2-03
% do_kho_du_kien: 5
Cho phương trình $x^2 - 2mx + m^2 - m + 2 = 0$ $(1)$, với $m$ là tham số.

a) Tìm $m$ để phương trình $(1)$ có hai nghiệm $x_1, x_2$ là độ dài hai cạnh góc vuông của một tam giác vuông có cạnh huyền bằng $6$.

b) Với giá trị $m$ tìm được, tính diện tích tam giác vuông đó.

\begin{solution}
a) Phương trình $(1)$ có $a = 1$, $b' = -m$, $c = m^2 - m + 2$, nên $\Delta' = m^2 - (m^2 - m + 2) = m - 2$. Phương trình $(1)$ có hai nghiệm (phân biệt hoặc kép) khi và chỉ khi $\Delta' \geq 0 \Leftrightarrow m \geq 2$.

a) Khi đó, theo định lý Viète: $x_1 + x_2 = 2m$ và $x_1 x_2 = m^2 - m + 2$.

a) $x_1, x_2$ là độ dài các cạnh nên phải dương, tức là $x_1 x_2 > 0$ (hai nghiệm cùng dấu) và $x_1 + x_2 > 0$ (cùng dương). Ta có $x_1 x_2 = \left(m - \dfrac{1}{2}\right)^2 + \dfrac{7}{4} > 0$ với mọi $m$, và $x_1 + x_2 = 2m > 0$ khi $m \geq 2$. Vậy với $m \geq 2$ hai nghiệm đều dương.

a) Theo định lý Pythagore, $x_1^2 + x_2^2 = 6^2 = 36$. Biến đổi theo tổng và tích: $x_1^2 + x_2^2 = (x_1 + x_2)^2 - 2x_1x_2 = 4m^2 - 2(m^2 - m + 2) = 2m^2 + 2m - 4$.

a) Điều kiện trở thành $2m^2 + 2m - 4 = 36 \Leftrightarrow m^2 + m - 20 = 0$; phương trình này có $\Delta = 1 + 80 = 81$, nên $m = \dfrac{-1 + 9}{2} = 4$ hoặc $m = \dfrac{-1 - 9}{2} = -5$.

a) Đối chiếu điều kiện $m \geq 2$: $m = 4$ thoả mãn; $m = -5$ bị loại (khi đó $\Delta' = -7 < 0$, phương trình $(1)$ vô nghiệm).

a) Thử lại với $m = 4$: $(1)$ là $x^2 - 8x + 14 = 0$, có $\Delta' = 16 - 14 = 2 > 0$, hai nghiệm $4 + \sqrt{2}$ và $4 - \sqrt{2}$ đều dương, và $(4 + \sqrt{2})^2 + (4 - \sqrt{2})^2 = 2(16 + 2) = 36$ (đúng).

b) Diện tích tam giác vuông bằng nửa tích hai cạnh góc vuông: $S = \dfrac{1}{2}x_1x_2 = \dfrac{1}{2}(4^2 - 4 + 2) = \dfrac{1}{2} \cdot 14 = 7$ (đơn vị diện tích).

Đáp số: a) $\boxed{m = 4}$; b) $\boxed{S = 7}$
\end{solution}
\end{problem}

% ---------------------------------------------------------------- phương trình khuyết: tìm chỗ sai trong lời giải
\begin{problem}[9DS-PT2-021]
% nguon: dummy (tự sinh để test)
% lop: 9
% kc_du_kien: 9DS-PT2-01, 8DS-PTNT-01
% do_kho_du_kien: 1
Hai bạn giải phương trình như sau:

a) Bạn An giải phương trình $3x^2 = 5x$: chia hai vế cho $x$ được $3x = 5$, nên phương trình có nghiệm duy nhất $x = \dfrac{5}{3}$.

b) Bạn Bình giải phương trình $2x^2 - 50 = 0$: biến đổi được $x^2 = 25$, nên $x = 5$.

Mỗi bạn giải đúng hay sai? Nếu sai, hãy chỉ ra chỗ sai và giải lại cho đúng.

\begin{solution}
a) Lời giải của An sai: An chia hai vế cho $x$ khi chưa biết $x \neq 0$, nên làm mất nghiệm $x = 0$.

a) Giải lại: chuyển vế được phương trình bậc hai khuyết hệ số tự do $3x^2 - 5x = 0$; đặt nhân tử chung: $x(3x - 5) = 0$.

a) Tích bằng $0$ khi một trong các thừa số bằng $0$: $x = 0$ hoặc $3x - 5 = 0 \Leftrightarrow x = \dfrac{5}{3}$.

b) Lời giải của Bình sai: số dương $25$ có hai căn bậc hai là $5$ và $-5$, nên phương trình $x^2 = 25$ có hai nghiệm; Bình bỏ sót nghiệm $x = -5$.

b) Giải lại: $2x^2 - 50 = 0 \Leftrightarrow 2x^2 = 50 \Leftrightarrow x^2 = 25 \Leftrightarrow x = 5$ hoặc $x = -5$.

Đáp số: a) $\boxed{\text{An sai; } x = 0;\ x = \dfrac{5}{3}}$; b) $\boxed{\text{Bình sai; } x = 5;\ x = -5}$
\end{solution}
\end{problem}

% ---------------------------------------------------------------- công thức nghiệm thu gọn với hệ số chứa căn; a, c trái dấu
\begin{problem}[9DS-PT2-022]
% nguon: dummy (tự sinh để test)
% lop: 9
% kc_du_kien: 9DS-PT2-02
% do_kho_du_kien: 2
a) Giải phương trình $x^2 - 2\sqrt{3}x - 6 = 0$.

b) Không giải phương trình, hãy giải thích vì sao phương trình $-3x^2 + 5x + 2026 = 0$ có hai nghiệm phân biệt.

\begin{solution}
a) Phương trình có $a = 1$, $b = -2\sqrt{3}$ nên $b' = -\sqrt{3}$, và $c = -6$. Dùng công thức nghiệm thu gọn: $\Delta' = b'^2 - ac = (-\sqrt{3})^2 - 1 \cdot (-6) = 3 + 6 = 9 > 0$, $\sqrt{\Delta'} = 3$.

a) Phương trình có hai nghiệm phân biệt: $x_1 = \dfrac{-b' + \sqrt{\Delta'}}{a} = \sqrt{3} + 3$ và $x_2 = \dfrac{-b' - \sqrt{\Delta'}}{a} = \sqrt{3} - 3$.

b) Phương trình có $a = -3$, $b = 5$, $c = 2026$, nên $ac = -3 \cdot 2026 = -6\,078 < 0$.

b) Biệt thức $\Delta = b^2 - 4ac$ có $b^2 = 25 \geq 0$ và $-4ac = 24\,312 > 0$, nên $\Delta = 25 + 24\,312 > 0$.

b) $\Delta > 0$ nên phương trình có hai nghiệm phân biệt. Tổng quát: mọi phương trình bậc hai có $a$ và $c$ trái dấu đều có hai nghiệm phân biệt.

Đáp số: a) $\boxed{x = 3 + \sqrt{3};\ x = \sqrt{3} - 3}$; b) $\boxed{ac < 0 \text{ nên } \Delta = b^2 - 4ac > 0}$
\end{solution}
\end{problem}

% ---------------------------------------------------------------- Viète: xét dấu các nghiệm không cần giải
\begin{problem}[9DS-PT2-023]
% nguon: dummy (tự sinh để test)
% lop: 9
% kc_du_kien: 9DS-PT2-03, 9DS-PT2-02
% do_kho_du_kien: 2
Không giải phương trình, hãy xét dấu các nghiệm (nếu có) của mỗi phương trình sau:

a) $3x^2 - 11x + 5 = 0$;

b) $2x^2 + 9x + 4 = 0$;

c) $5x^2 + 3x - 7 = 0$.

\begin{solution}
a) $\Delta = (-11)^2 - 4 \cdot 3 \cdot 5 = 121 - 60 = 61 > 0$ nên phương trình có hai nghiệm phân biệt $x_1, x_2$.

a) Theo định lý Viète: $x_1 + x_2 = \dfrac{11}{3} > 0$ và $x_1 x_2 = \dfrac{5}{3} > 0$. Tích dương nên hai nghiệm cùng dấu; tổng dương nên cả hai nghiệm đều dương.

b) $\Delta = 9^2 - 4 \cdot 2 \cdot 4 = 81 - 32 = 49 > 0$ nên phương trình có hai nghiệm phân biệt $x_1, x_2$.

b) Theo định lý Viète: $x_1 + x_2 = -\dfrac{9}{2} < 0$ và $x_1 x_2 = \dfrac{4}{2} = 2 > 0$. Tích dương nên hai nghiệm cùng dấu; tổng âm nên cả hai nghiệm đều âm.

c) $a = 5$ và $c = -7$ trái dấu nên $ac < 0$, suy ra $\Delta = b^2 - 4ac > 0$: phương trình có hai nghiệm phân biệt $x_1, x_2$.

c) Theo định lý Viète: $x_1 x_2 = -\dfrac{7}{5} < 0$ nên hai nghiệm trái dấu.

Đáp số: a) $\boxed{\text{hai nghiệm dương phân biệt}}$; b) $\boxed{\text{hai nghiệm âm phân biệt}}$; c) $\boxed{\text{hai nghiệm trái dấu}}$
\end{solution}
\end{problem}

% ---------------------------------------------------------------- Viète: biểu thức không đối xứng (hạ bậc), tổng căn các nghiệm
\begin{problem}[9DS-PT2-024]
% nguon: dummy (tự sinh để test)
% lop: 9
% kc_du_kien: 9DS-PT2-03, 8DS-HDT-04, 9DS-PT2-02
% do_kho_du_kien: 5
Cho phương trình $x^2 - 3x + 1 = 0$.

a) Chứng tỏ rằng phương trình có hai nghiệm dương phân biệt $x_1, x_2$.

b) Không giải phương trình, tính giá trị của biểu thức $A = x_1^3 + 8x_2$.

c) Không giải phương trình, tính $B = \sqrt{x_1} + \sqrt{x_2}$ và $C = x_1\sqrt{x_1} + x_2\sqrt{x_2}$.

\begin{solution}
a) $\Delta = (-3)^2 - 4 \cdot 1 \cdot 1 = 5 > 0$ nên phương trình có hai nghiệm phân biệt $x_1, x_2$. Theo định lý Viète: $x_1 + x_2 = 3 > 0$ và $x_1 x_2 = 1 > 0$, nên hai nghiệm cùng dấu và đều dương.

b) Vì $x_1$ là nghiệm nên $x_1^2 - 3x_1 + 1 = 0$, tức là $x_1^2 = 3x_1 - 1$ (dùng để hạ bậc).

b) Nhân hai vế với $x_1$ rồi thay $x_1^2 = 3x_1 - 1$ lần nữa: $x_1^3 = 3x_1^2 - x_1 = 3(3x_1 - 1) - x_1 = 8x_1 - 3$.

b) Do đó $A = (8x_1 - 3) + 8x_2 = 8(x_1 + x_2) - 3 = 8 \cdot 3 - 3 = 21$ (theo định lý Viète).

c) Vì $x_1 > 0$, $x_2 > 0$ nên $\sqrt{x_1}$, $\sqrt{x_2}$ xác định và $B > 0$. Ta có $B^2 = x_1 + x_2 + 2\sqrt{x_1}\sqrt{x_2} = (x_1 + x_2) + 2\sqrt{x_1 x_2} = 3 + 2\sqrt{1} = 5$, nên $B = \sqrt{5}$.

c) Đặt $u = \sqrt{x_1}$, $v = \sqrt{x_2}$ thì $u + v = \sqrt{5}$, $uv = \sqrt{x_1 x_2} = 1$ và $C = u^3 + v^3$. Theo hằng đẳng thức tổng hai lập phương: $u^3 + v^3 = (u + v)(u^2 - uv + v^2) = (u + v)[(u + v)^2 - 3uv]$.

c) Thay số: $C = \sqrt{5} \cdot (5 - 3 \cdot 1) = 2\sqrt{5}$.

Đáp số: a) $\boxed{\Delta = 5 > 0,\ x_1 + x_2 = 3 > 0,\ x_1 x_2 = 1 > 0}$; b) $\boxed{A = 21}$; c) $\boxed{B = \sqrt{5},\ C = 2\sqrt{5}}$
\end{solution}
\end{problem}

% ---------------------------------------------------------------- nhẩm nghiệm để phân tích tam thức thành nhân tử, rút gọn phân thức
\begin{problem}[9DS-PT2-025]
% nguon: dummy (tự sinh để test)
% lop: 9
% kc_du_kien: 9DS-PT2-04, 8DS-PS-02, 9DS-PT2-03
% do_kho_du_kien: 3
a) Nhẩm nghiệm của các phương trình $3x^2 - 8x + 5 = 0$ và $2x^2 + 3x - 5 = 0$.

b) Dùng kết quả câu a), phân tích các đa thức $3x^2 - 8x + 5$ và $2x^2 + 3x - 5$ thành nhân tử.

c) Rút gọn phân thức $P = \dfrac{3x^2 - 8x + 5}{2x^2 + 3x - 5}$.

\begin{solution}
a) Phương trình $3x^2 - 8x + 5 = 0$ có $a + b + c = 3 - 8 + 5 = 0$ nên có hai nghiệm $x_1 = 1$ và $x_2 = \dfrac{c}{a} = \dfrac{5}{3}$.

a) Phương trình $2x^2 + 3x - 5 = 0$ có $a + b + c = 2 + 3 - 5 = 0$ nên có hai nghiệm $x_1 = 1$ và $x_2 = \dfrac{c}{a} = -\dfrac{5}{2}$.

b) Nếu $ax^2 + bx + c = 0$ có hai nghiệm $x_1, x_2$ thì theo định lý Viète $b = -a(x_1 + x_2)$ và $c = ax_1x_2$, nên $ax^2 + bx + c = a[x^2 - (x_1 + x_2)x + x_1x_2] = a(x - x_1)(x - x_2)$.

b) Áp dụng với các nghiệm ở câu a): $3x^2 - 8x + 5 = 3(x - 1)\left(x - \dfrac{5}{3}\right) = (x - 1)(3x - 5)$ và $2x^2 + 3x - 5 = 2(x - 1)\left(x + \dfrac{5}{2}\right) = (x - 1)(2x + 5)$.

c) Phân thức xác định khi mẫu khác $0$: $(x - 1)(2x + 5) \neq 0$, tức là $x \neq 1$ và $x \neq -\dfrac{5}{2}$.

c) Chia cả tử và mẫu cho nhân tử chung $x - 1 \neq 0$: $P = \dfrac{(x - 1)(3x - 5)}{(x - 1)(2x + 5)} = \dfrac{3x - 5}{2x + 5}$.

Đáp số: a) $\boxed{x = 1,\ x = \dfrac{5}{3} \text{ và } x = 1,\ x = -\dfrac{5}{2}}$; b) $\boxed{(x - 1)(3x - 5) \text{ và } (x - 1)(2x + 5)}$; c) $\boxed{P = \dfrac{3x - 5}{2x + 5} \text{ với } x \neq 1,\ x \neq -\dfrac{5}{2}}$
\end{solution}
\end{problem}

% ---------------------------------------------------------------- nhẩm nghiệm với hệ số chứa căn, trục căn thức nghiệm thứ hai
\begin{problem}[9DS-PT2-026]
% nguon: dummy (tự sinh để test)
% lop: 9
% kc_du_kien: 9DS-PT2-04, 9DS-CAN-04
% do_kho_du_kien: 3
Nhẩm nghiệm của các phương trình sau (viết các nghiệm dưới dạng không còn căn thức ở mẫu):

a) $(2 - \sqrt{3})x^2 + 2\sqrt{3}x - (2 + \sqrt{3}) = 0$;

b) $(\sqrt{3} + 1)x^2 + 2\sqrt{3}x + \sqrt{3} - 1 = 0$.

\begin{solution}
a) Phương trình có $a = 2 - \sqrt{3} \neq 0$, $b = 2\sqrt{3}$, $c = -(2 + \sqrt{3})$ và $a + b + c = 2 - \sqrt{3} + 2\sqrt{3} - 2 - \sqrt{3} = 0$.

a) Do đó phương trình có hai nghiệm $x_1 = 1$ và $x_2 = \dfrac{c}{a} = -\dfrac{2 + \sqrt{3}}{2 - \sqrt{3}}$.

a) Trục căn thức ở mẫu: nhân cả tử và mẫu với $2 + \sqrt{3}$, dùng $(2 - \sqrt{3})(2 + \sqrt{3}) = 4 - 3 = 1$: $x_2 = -(2 + \sqrt{3})^2 = -(4 + 4\sqrt{3} + 3) = -7 - 4\sqrt{3}$.

b) Phương trình có $a = \sqrt{3} + 1 \neq 0$, $b = 2\sqrt{3}$, $c = \sqrt{3} - 1$ và $a - b + c = \sqrt{3} + 1 - 2\sqrt{3} + \sqrt{3} - 1 = 0$.

b) Do đó phương trình có hai nghiệm $x_1 = -1$ và $x_2 = -\dfrac{c}{a} = -\dfrac{\sqrt{3} - 1}{\sqrt{3} + 1}$.

b) Trục căn thức ở mẫu: nhân cả tử và mẫu với $\sqrt{3} - 1$, dùng $(\sqrt{3} + 1)(\sqrt{3} - 1) = 3 - 1 = 2$: $x_2 = -\dfrac{(\sqrt{3} - 1)^2}{2} = -\dfrac{4 - 2\sqrt{3}}{2} = \sqrt{3} - 2$.

Đáp số: a) $\boxed{x_1 = 1;\ x_2 = -7 - 4\sqrt{3}}$; b) $\boxed{x_1 = -1;\ x_2 = \sqrt{3} - 2}$
\end{solution}
\end{problem}

% ---------------------------------------------------------------- nhẩm nghiệm với hệ số bằng chữ, điều kiện nghiệm kép, hai cách
\begin{problem}[9DS-PT2-027]
% nguon: dummy (tự sinh để test)
% lop: 9
% kc_du_kien: 9DS-PT2-04, 9DS-PT2-02, 8DS-HDT-01
% do_kho_du_kien: 4
Cho phương trình ẩn $x$: $(a - b)x^2 + (b - c)x + c - a = 0$ $(1)$, trong đó $a, b, c$ là các số thực và $a \neq b$.

a) Chứng minh rằng phương trình $(1)$ luôn có nghiệm và tìm các nghiệm của nó.

b) Tìm hệ thức giữa $a, b, c$ để phương trình $(1)$ có nghiệm kép.

\begin{solution}[cách 1: nhẩm nghiệm]
a) Vì $a \neq b$ nên hệ số của $x^2$ là $a - b \neq 0$, do đó $(1)$ là phương trình bậc hai. Tổng các hệ số: $(a - b) + (b - c) + (c - a) = 0$.

a) Theo quy tắc nhẩm nghiệm (tổng các hệ số bằng $0$), $(1)$ luôn có hai nghiệm $x_1 = 1$ và $x_2 = \dfrac{c - a}{a - b}$ (có thể trùng nhau).

b) Phương trình bậc hai $(1)$ có nghiệm kép khi và chỉ khi hai nghiệm này trùng nhau: $x_2 = x_1 \Leftrightarrow \dfrac{c - a}{a - b} = 1$.

b) Vì $a - b \neq 0$, nhân hai vế với $a - b$: $c - a = a - b \Leftrightarrow b + c = 2a$. Khi đó nghiệm kép là $x = 1$.

Đáp số: a) $\boxed{x_1 = 1;\ x_2 = \dfrac{c - a}{a - b}}$; b) $\boxed{b + c = 2a}$
\end{solution}

\begin{solution}[cách 2: dùng biệt thức]
a) Vì $a - b \neq 0$ nên $(1)$ là phương trình bậc hai, có $\Delta = (b - c)^2 - 4(a - b)(c - a)$.

a) Đặt $p = a - b$, $q = c - a$ thì $p + q = c - b = -(b - c)$. Theo hằng đẳng thức bình phương của một tổng, một hiệu: $\Delta = (p + q)^2 - 4pq = (p - q)^2 = (2a - b - c)^2 \geq 0$, nên $(1)$ luôn có nghiệm.

a) Theo công thức nghiệm, với $-(b - c) = p + q$ và $\sqrt{\Delta} = |p - q|$ (vì $\pm|p - q|$ và $\pm(p - q)$ cho cùng hai giá trị): $x = \dfrac{(p + q) \pm (p - q)}{2p}$, cho hai nghiệm $\dfrac{2p}{2p} = 1$ và $\dfrac{2q}{2p} = \dfrac{c - a}{a - b}$.

b) Phương trình có nghiệm kép khi và chỉ khi $\Delta = 0 \Leftrightarrow (2a - b - c)^2 = 0 \Leftrightarrow b + c = 2a$.

Đáp số: a) $\boxed{x_1 = 1;\ x_2 = \dfrac{c - a}{a - b}}$; b) $\boxed{b + c = 2a}$
\end{solution}
\end{problem}

% ---------------------------------------------------------------- lập phương trình bậc hai biết hai nghiệm, hai cách
\begin{problem}[9DS-PT2-028]
% nguon: dummy (tự sinh để test)
% lop: 9
% kc_du_kien: 9DS-PT2-05, 8DS-HDT-02
% do_kho_du_kien: 2
Lập một phương trình bậc hai (với các hệ số nguyên) có hai nghiệm là:

a) $2 + \sqrt{3}$ và $2 - \sqrt{3}$;

b) $\dfrac{1}{3}$ và $-\dfrac{5}{2}$.

\begin{solution}[cách 1: dùng tổng và tích]
a) Tổng hai số: $S = (2 + \sqrt{3}) + (2 - \sqrt{3}) = 4$; tích hai số, theo hằng đẳng thức hiệu hai bình phương: $P = (2 + \sqrt{3})(2 - \sqrt{3}) = 4 - 3 = 1$.

a) Hai số có tổng $S$ và tích $P$ là hai nghiệm của phương trình $X^2 - SX + P = 0$ (điều kiện $S^2 \geq 4P$ thoả mãn vì $16 \geq 4$), nên phương trình cần lập là $x^2 - 4x + 1 = 0$.

b) Tổng hai số: $S = \dfrac{1}{3} - \dfrac{5}{2} = \dfrac{2 - 15}{6} = -\dfrac{13}{6}$; tích: $P = \dfrac{1}{3} \cdot \left(-\dfrac{5}{2}\right) = -\dfrac{5}{6}$ (điều kiện $S^2 \geq 4P$ hiển nhiên đúng vì $P < 0$).

b) Hai số đó là nghiệm của $x^2 + \dfrac{13}{6}x - \dfrac{5}{6} = 0$; nhân hai vế với $6$ để được hệ số nguyên: $6x^2 + 13x - 5 = 0$.

Đáp số: a) $\boxed{x^2 - 4x + 1 = 0}$; b) $\boxed{6x^2 + 13x - 5 = 0}$
\end{solution}

\begin{solution}[cách 2: viết phương trình dạng tích]
a) Phương trình $(x - x_1)(x - x_2) = 0$ có đúng hai nghiệm $x_1, x_2$ và khai triển ra là phương trình bậc hai. Với $x_1 = 2 + \sqrt{3}$, $x_2 = 2 - \sqrt{3}$: $(x - 2 - \sqrt{3})(x - 2 + \sqrt{3}) = 0$.

a) Dùng hằng đẳng thức hiệu hai bình phương với $A = x - 2$, $B = \sqrt{3}$: $(x - 2)^2 - 3 = 0 \Leftrightarrow x^2 - 4x + 1 = 0$.

b) Phương trình $\left(x - \dfrac{1}{3}\right)\left(x + \dfrac{5}{2}\right) = 0$ có hai nghiệm $\dfrac{1}{3}$ và $-\dfrac{5}{2}$; nhân hai vế với $6 = 3 \cdot 2$: $(3x - 1)(2x + 5) = 0$.

b) Khai triển vế trái: $6x^2 + 15x - 2x - 5 = 0 \Leftrightarrow 6x^2 + 13x - 5 = 0$.

Đáp số: a) $\boxed{x^2 - 4x + 1 = 0}$; b) $\boxed{6x^2 + 13x - 5 = 0}$
\end{solution}
\end{problem}

% ---------------------------------------------------------------- hệ đối xứng: đặt tổng và tích, tìm hai số
\begin{problem}[9DS-PT2-029]
% nguon: dummy (tự sinh để test)
% lop: 9
% kc_du_kien: 9DS-PT2-05, 9DS-PT2-02, 8DS-HDT-01
% do_kho_du_kien: 4
Giải hệ phương trình $\begin{cases} x^2 + y^2 + x + y = 18 \\ xy = 6 \end{cases}$

\begin{solution}
Đặt $S = x + y$ và $P = xy$. Theo hằng đẳng thức bình phương của một tổng: $x^2 + y^2 = (x + y)^2 - 2xy = S^2 - 2P$.

Hệ trở thành $\begin{cases} S^2 - 2P + S = 18 \\ P = 6 \end{cases}$; thay $P = 6$ vào phương trình đầu được $S^2 + S - 30 = 0$.

Giải phương trình bậc hai ẩn $S$: $\Delta = 1^2 + 4 \cdot 30 = 121$, $\sqrt{\Delta} = 11$, nên $S = \dfrac{-1 + 11}{2} = 5$ hoặc $S = \dfrac{-1 - 11}{2} = -6$.

Trường hợp $S = 5$, $P = 6$: điều kiện $S^2 \geq 4P$ thoả mãn ($25 \geq 24$), nên $x, y$ là hai nghiệm của $X^2 - 5X + 6 = 0$. Phương trình này có $\Delta = 25 - 24 = 1$, nghiệm $X = \dfrac{5 + 1}{2} = 3$ và $X = \dfrac{5 - 1}{2} = 2$; suy ra $(x; y) = (2; 3)$ hoặc $(x; y) = (3; 2)$.

Trường hợp $S = -6$, $P = 6$: điều kiện $S^2 \geq 4P$ thoả mãn ($36 \geq 24$), nên $x, y$ là hai nghiệm của $X^2 + 6X + 6 = 0$. Dùng công thức nghiệm thu gọn: $\Delta' = 3^2 - 6 = 3$, nghiệm $X = -3 + \sqrt{3}$ và $X = -3 - \sqrt{3}$; suy ra $(x; y) = (-3 + \sqrt{3}; -3 - \sqrt{3})$ hoặc $(x; y) = (-3 - \sqrt{3}; -3 + \sqrt{3})$.

Thử lại: với $(2; 3)$: $4 + 9 + 2 + 3 = 18$ và $2 \cdot 3 = 6$ (đúng); với $x = -3 + \sqrt{3}$, $y = -3 - \sqrt{3}$: $xy = 9 - 3 = 6$, $x + y = -6$, $x^2 + y^2 = 36 - 12 = 24$, nên $24 + (-6) = 18$ (đúng); hai cặp còn lại có được do đổi chỗ $x$ và $y$.

Đáp số: $\boxed{(x; y) = (2; 3),\ (3; 2),\ (-3 + \sqrt{3}; -3 - \sqrt{3}),\ (-3 - \sqrt{3}; -3 + \sqrt{3})}$
\end{solution}
\end{problem}

% ---------------------------------------------------------------- bài toán thực tế: gấp hộp không nắp từ tấm tôn
\begin{problem}[9DS-PT2-030]
% nguon: dummy (tự sinh để test)
% lop: 9
% kc_du_kien: 9DS-PT2-07, 9DS-PT2-01
% do_kho_du_kien: 3
Từ một tấm tôn hình vuông, người ta cắt bỏ ở bốn góc bốn hình vuông bằng nhau có cạnh $5$ cm, rồi gấp phần còn lại thành một chiếc hộp không nắp dạng hình hộp chữ nhật có thể tích $1\,125 \text{ cm}^3$. Tính độ dài cạnh của tấm tôn ban đầu.

\begin{solution}
Gọi độ dài cạnh của tấm tôn ban đầu là $x$ (cm). Để cắt được bốn góc và đáy hộp có cạnh dương, cần $x - 2 \cdot 5 > 0$, nên điều kiện là $x > 10$.

Sau khi cắt và gấp, đáy hộp là hình vuông cạnh $x - 10$ (cm) và chiều cao hộp bằng cạnh hình vuông bị cắt, tức $5$ cm.

Thể tích hộp bằng diện tích đáy nhân chiều cao, nên ta có phương trình $5(x - 10)^2 = 1\,125$ (tức là phương trình bậc hai $x^2 - 20x - 125 = 0$).

Chia hai vế cho $5$: $(x - 10)^2 = 225$. Vì $225 > 0$ có hai căn bậc hai là $15$ và $-15$, nên $x - 10 = 15$ hoặc $x - 10 = -15$, tức là $x = 25$ hoặc $x = -5$.

Đối chiếu điều kiện $x > 10$: $x = 25$ thoả mãn, $x = -5$ bị loại.

Thử lại: đáy hộp có cạnh $25 - 10 = 15$ cm, thể tích hộp $15^2 \cdot 5 = 1\,125 \text{ cm}^3$ (đúng).

Đáp số: $\boxed{25 \text{ cm}}$
\end{solution}
\end{problem}

% ---------------------------------------------------------------- bài toán năng suất: hai vòi nước cùng chảy
\begin{problem}[9DS-PT2-031]
% nguon: dummy (tự sinh để test)
% lop: 9
% kc_du_kien: 9DS-PT2-07, 9DS-PT2-02, 8DS-PS-03
% do_kho_du_kien: 4
Hai vòi nước cùng chảy vào một bể cạn thì sau $4$ giờ $48$ phút đầy bể. Nếu chảy riêng thì vòi thứ nhất chảy đầy bể nhanh hơn vòi thứ hai $4$ giờ. Hỏi nếu chảy riêng thì mỗi vòi chảy đầy bể trong bao lâu?

\begin{solution}
Đổi đơn vị: $4$ giờ $48$ phút $= 4 + \dfrac{48}{60} = \dfrac{24}{5}$ giờ.

Gọi thời gian vòi thứ nhất chảy riêng đầy bể là $x$ (giờ), điều kiện $x > \dfrac{24}{5}$ (mỗi vòi chảy riêng phải lâu hơn khi hai vòi cùng chảy). Khi đó thời gian vòi thứ hai chảy riêng đầy bể là $x + 4$ (giờ).

Trong $1$ giờ, vòi thứ nhất chảy được $\dfrac{1}{x}$ bể, vòi thứ hai chảy được $\dfrac{1}{x + 4}$ bể, cả hai vòi chảy được $1 : \dfrac{24}{5} = \dfrac{5}{24}$ bể, nên ta có phương trình $\dfrac{1}{x} + \dfrac{1}{x + 4} = \dfrac{5}{24}$.

Với $x > \dfrac{24}{5}$ thì $x(x + 4) \neq 0$; quy đồng với mẫu chung $24x(x + 4)$ rồi khử mẫu: $24(x + 4) + 24x = 5x(x + 4) \Leftrightarrow 48x + 96 = 5x^2 + 20x \Leftrightarrow 5x^2 - 28x - 96 = 0$.

Dùng công thức nghiệm thu gọn với $b' = -14$: $\Delta' = (-14)^2 - 5 \cdot (-96) = 196 + 480 = 676$, $\sqrt{\Delta'} = 26$, nên $x_1 = \dfrac{14 + 26}{5} = 8$ và $x_2 = \dfrac{14 - 26}{5} = -\dfrac{12}{5}$.

Đối chiếu điều kiện $x > \dfrac{24}{5}$: $x_1 = 8$ thoả mãn, $x_2 = -\dfrac{12}{5}$ bị loại. Vậy vòi thứ nhất chảy riêng đầy bể trong $8$ giờ, vòi thứ hai trong $8 + 4 = 12$ giờ.

Thử lại: $\dfrac{1}{8} + \dfrac{1}{12} = \dfrac{3 + 2}{24} = \dfrac{5}{24}$ (đúng).

Đáp số: $\boxed{\text{vòi thứ nhất: } 8 \text{ giờ};\ \text{vòi thứ hai: } 12 \text{ giờ}}$
\end{solution}
\end{problem}

\end{document}
